\documentclass[11pt]{article}
\usepackage[a4paper,margin=25mm]{geometry}
\usepackage{amsmath,amssymb,amsthm}
\usepackage[hidelinks]{hyperref}
\newtheorem{theorem}{Theorem}
\newtheorem{lemma}{Lemma}
\title{A Nine-Vertex Counterexample to Hypergraph Reconstruction:\\Disproof of Conjecture 00000001683}
\author{GitHub: gaochengzhecpu}
\date{}
\begin{document}
\maketitle
\begin{abstract}
Two explicit simple $3$-uniform hypergraphs on nine vertices have equal
vertex-deleted decks, including multiplicities, but are not isomorphic.
Neither is a Steiner triple system. Thus they refute both the assertion
of reconstruction for every order at least eight and the assertion that
every counterexample must consist of Steiner triple systems. The example
is a classical construction of Kocay; the contribution here is an explicit,
self-contained certificate and its kernel-checked Lean verification.
\end{abstract}

\paragraph{Reading of the source.}
Both source languages assert reconstruction for $n\geq8$ and state that
any counterexample must be a twin pair of Steiner triple systems. We use
the ordinary vertex deck: the multiset of isomorphism classes of the
induced hypergraphs obtained by deleting one vertex. A simple
$3$-uniform hypergraph on $V$ is a set of three-element subsets of $V$.
An isomorphism is a bijection of vertices carrying the edge set exactly
onto the other edge set. A Steiner triple system requires every pair of
distinct vertices to belong to exactly one edge. In particular, it is a
$2$-$(v,3,1)$ design; the source's phrase ``3-design parameters'' is not
needed for the argument.

\section*{The explicit pair}
Write $abc$ for the unordered triple $\{a,b,c\}$. Let
$V=\{0,1,\ldots,8\}$ and let the following twenty triples form $B$:
\[
\begin{aligned}
B=\{&123,128,135,136,146,147,178,234,246,247,\\
    &257,258,345,357,358,368,456,468,567,678\}.
\end{aligned}
\]
Define
\[
\begin{aligned}
A_X&=\{012,016,025,034,038,047,056,078\},\\
A_Y&=\{014,018,023,027,036,045,058,067\},\\
E(X)&=B\cup A_X,\qquad E(Y)=B\cup A_Y.
\end{aligned}
\]
All listed triples have three distinct vertices, and within either edge
set there are no repetitions. Thus both hypergraphs are simple,
$3$-uniform, and have $28$ edges.

\begin{theorem}
The hypergraphs $X,Y$ have the same vertex deck but are not isomorphic.
Neither $X$ nor $Y$ is a Steiner triple system. Consequently both clauses
of the source described above are false.
\end{theorem}

\begin{proof}
We prove the three assertions in turn.

\emph{Equal decks.}
The table gives permutations $f$ of $V$. The entries in the last nine
columns are $f(0),\ldots,f(8)$. Each row is used for all deleted vertices
in its first column.
\[
\begin{array}{c|rrrrrrrrr}
\text{deleted }v&0&1&2&3&4&5&6&7&8\\\hline
0   &0&1&2&3&4&5&6&7&8\\
1,5 &0&1&8&7&6&5&4&3&2\\
2,6 &0&7&2&5&8&3&6&1&4\\
3,7 &0&5&8&3&6&1&4&7&2\\
4,8 &0&7&6&5&4&3&2&1&8
\end{array}
\]
Each row is an involution and fixes the relevant deleted vertex $v$.
Substitution into the displayed edge sets gives the following exact
equalities:
\[
\begin{array}{c|c|c}
v&\#E(X-v)=\#E(Y-v)& f(E(X-v))=E(Y-v)\\\hline
0&20&\text{yes}\\
1,2,7,8&19&\text{yes, using the row for each }v\\
3,4,5,6&18&\text{yes, using the row for each }v
\end{array}
\]
These are finite equalities of sets of triples: apply the row permutation
to every listed triple that avoids $v$, sort its entries, and compare
with the other list after deleting triples containing $v$.
They prove $X-v\cong Y-v$ for every $v$. The identity matching of the
nine deleted vertices preserves all deck multiplicities.

\emph{Non-isomorphism.}
Let $d_H(a)$ count edges containing $a$, and let $c_H(a,b)$ count edges
containing the distinct vertices $a,b$. An isomorphism preserves both
quantities. The degree sequence in vertex order is, for both graphs,
\[
\bigl(d_H(0),\ldots,d_H(8)\bigr)=(8,9,9,10,10,10,10,9,9).
\]
Thus any isomorphism $e:X\longrightarrow Y$ must fix $0$, preserve
$L=\{1,2,7,8\}$, and preserve $D=\{3,4,5,6\}$.
Since $c_X(1,2)=3$, the ordered pair $(e(1),e(2))$ must be one of
\[
(1,8),\quad(2,7),\quad(7,2),\quad(8,1).
\]
Indeed the only unordered pairs in $L$ of $Y$ with codegree three are
$\{1,8\}$ and $\{2,7\}$. For the next step, the codegrees in $Y$ between
$D$ and $L$ are
\[
\begin{array}{c|rrrr}
c_Y&1&2&7&8\\\hline
3&3&3&1&2\\
4&3&3&2&1\\
5&1&2&3&3\\
6&2&1&2&3
\end{array}
\]
In $X$, we have
\[
(c_X(3,1),c_X(3,2))=(3,2),\qquad
(c_X(4,1),c_X(4,2))=(2,3).
\]
The table therefore forces the additional image indicated in each row
below. The last two columns give a contradiction to edge preservation.
\[
\begin{array}{cc|c|cc|cc}
e(1)&e(2)&\text{forced image}&T&e(T)&T\in E(X)&e(T)\in E(Y)\\\hline
1&8&e(3)=3&123&138&\text{yes}&\text{no}\\
2&7&e(4)=5&124&257&\text{no}&\text{yes}\\
7&2&e(4)=4&124&247&\text{no}&\text{yes}\\
8&1&e(3)=6&123&168&\text{yes}&\text{no}
\end{array}
\]
This exhausts all vertex bijections that could be isomorphisms.

\emph{Failure of the Steiner condition.}
In $X$, the pair $\{0,1\}$ belongs to the two edges $012,016$; in $Y$,
it belongs to $014,018$. In each case its codegree is two, whereas a
Steiner triple system requires codegree one. Since $9\geq8$, the pair
refutes reconstruction in the claimed range, and it also refutes the
claim that counterexamples must be Steiner systems.
\end{proof}

\paragraph{The congruence clause.}
The order nine satisfies $9\equiv3\pmod6$. We do not dispute the usual
necessary condition $v\equiv1,3\pmod6$ for a Steiner triple system.
The counterexample invalidates the asserted restriction to Steiner
systems, rather than that necessary condition itself. No meaning of
``twin'' can restore the assertion when neither member is Steiner.

\section*{Attribution}
This is the $n=3$ case, with eight ordinary vertices and the additional
vertex $0$, of Kocay's classical family \cite{kocay}. Cooper and Okur
\cite{cooperokur}, Definitions 9--13 and Theorem 2.1, describe the
construction. We give all edges and verify this finite pair directly;
we do not claim a new nonreconstruction construction. No theorem about
the infinite family or characteristic polynomials is used in our proof.

\section*{Formal verification and scope}
The Lean structure \texttt{Hypergraph3} is a Boolean relation on triples
with proved symmetry under adjacent swaps and vanishing on repeated
vertices. These properties express a simple $3$-uniform hypergraph.
The concrete objects \texttt{X} and \texttt{Y} use a bit indexed by
$81a+9b+c$ to represent the edge relation efficiently. The finite sets
\texttt{edgesX} and \texttt{edgesY} contain exactly the unordered triples
printed above. Their uniformity and cardinality $28$ are proved.
The bridge theorems
\begin{center}
\texttt{edgeX\_iff\_mem},\qquad\texttt{edgeY\_iff\_mem}
\end{center}
prove, for every $a,b,c\in\operatorname{Fin}(9)$, that the encoded edge
relation holds if and only if $\{a,b,c\}$ belongs to the corresponding
displayed edge set. Thus the bit encoding does not replace the
mathematical hypergraph with an unrelated numerical assertion.

\texttt{IsIso} quantifies over every triple and preserves edges in both
directions under an arbitrary vertex equivalence. \texttt{Hypomorphic}
contains an outer bijection matching the deleted vertices and an inner
vertex equivalence for each card; its restriction to the complements is
an isomorphism. Conversely, any card isomorphism extends to such an
equivalence by sending the omitted vertex to the omitted vertex. This
is exactly equality of decks with multiplicity, not a weaker comparison
of sets of card types. The finite equivalences and all card conditions
are checked in \texttt{same\_deck}.

\texttt{codegree} counts the actual third vertices completing a pair;
\texttt{doubleDegree} counts ordered pairs completing a vertex and is
twice its usual degree. Their preservation under arbitrary
\texttt{IsIso} is proved by equivalences of the counting types.
The theorem \texttt{not\_isomorphic} excludes every vertex equivalence,
using exhaustive finite cases pruned by these proved invariants and
by individual edge contradictions.

The proposition \texttt{ReconstructionClaim} quantifies over every
$n\geq8$ and all simple $3$-uniform hypergraphs on $\operatorname{Fin}(n)$.
The final theorem
\begin{center}\texttt{conjecture1683\_false : $\neg$ ReconstructionClaim}
\end{center}
is its negation. The theorems \texttt{X\_not\_steiner} and
\texttt{Y\_not\_steiner} use the usual pair-codegree definition of a
Steiner triple system. The theorem \texttt{steiner\_clause\_false}
negates the universal claim that hypomorphic nonisomorphic pairs of
order at least eight must both be Steiner systems. No statement about
the arithmetic existence conditions for designs is formalised.

All finite checks use kernel evaluation with \texttt{decide}.
The development contains no placeholder proofs, native decision
procedures, or custom axioms. Only three standard axioms occur in the audit:
\begin{center}
\texttt{propext},\quad \texttt{Classical.choice},\quad \texttt{Quot.sound}.
\end{center}

\section*{Source and reproduction}
The exact bilingual source is included as \texttt{SOURCE.md}, from
\begin{center}\small\texttt{conjectures/00000001683.md}.\end{center}
The source snapshot is pinned to upstream commit
\begin{center}\small\texttt{b355e42906e6eaf52bc7e9d005e0ef02d8aa7400}.\end{center}
The project uses Lean 4.19.0 and Mathlib commit
\begin{center}\small\texttt{c44e0c8ee63ca166450922a373c7409c5d26b00b}.\end{center}
From \texttt{lean/}, run \texttt{lake build}, then
\begin{center}\small\texttt{lake env lean -DwarningAsError=true Main.lean}.
\end{center}
The supplied manifest pins the dependencies. Compile this document with
Tectonic or a compatible LaTeX installation. No external numerical
program is required to check the formal proof.

\begin{thebibliography}{9}
\raggedright
\bibitem{kocay}
W.~L. Kocay, \emph{A family of nonreconstructible hypergraphs},
Journal of Combinatorial Theory, Series B \textbf{42} (1987), 46--63.
\bibitem{cooperokur}
J.~Cooper and U.~Okur, \emph{Polynomial Reconstruction Problem for
Hypergraphs}, arXiv:2312.16152v2.
\href{https://arxiv.org/html/2312.16152v2}{Public construction and attribution}.
\end{thebibliography}
\end{document}
